\documentclass[12pt,a4paper]{article}

% ---- Encoding & Font ----
\usepackage{iftex}
\ifPDFTeX
  \usepackage[utf8]{inputenc}
  \usepackage[T1]{fontenc}
\else
  \usepackage{fontspec}  % XeTeX/LuaTeX (tectonic): Unicode nativo (·, —, ¿, ª, §)
\fi
\usepackage[spanish]{babel}
\usepackage{csquotes}

% ---- Page layout ----
\usepackage[top=2.5cm, bottom=2.5cm, left=2.5cm, right=2.5cm]{geometry}
\usepackage{setspace}
\onehalfspacing

% ---- Math ----
\usepackage{amsmath, amssymb, amsthm}
\usepackage{mathtools}
\usepackage{physics}

% ---- Graphics ----
\usepackage{graphicx}
\usepackage{tikz}
\usepackage{float}
\usepackage{caption}

% ---- Tables ----
\usepackage{array}
\usepackage{booktabs}
\usepackage{tabularx}
\usepackage{colortbl}

% ---- Colours ----
\usepackage{xcolor}
\definecolor{notebg}{HTML}{FFF3CD}
\definecolor{noteborder}{HTML}{FFC107}
\definecolor{boxred}{HTML}{F8D7DA}
\definecolor{boxredborder}{HTML}{F5C6CB}

% ---- Boxes ----
\usepackage{mdframed}
\usepackage{tcolorbox}
\tcbuselibrary{most}

\newtcolorbox{keybox}[1][]{
  colback=blue!5,
  colframe=blue!70!black,
  arc=4pt,
  boxrule=1pt,
  fonttitle=\bfseries,
  title={#1}
}

\newtcolorbox{warnbox}[1][]{
  colback=boxred,
  colframe=boxredborder,
  coltitle=black!75,
  arc=4pt,
  boxrule=1pt,
  fonttitle=\bfseries,
  title={#1}
}

\newtcolorbox{notebox}[1][]{
  colback=notebg,
  colframe=noteborder,
  coltitle=black!75,
  arc=4pt,
  boxrule=1pt,
  fonttitle=\bfseries,
  title={#1}
}

% ---- Hyperlinks & TOC ----
\usepackage[hidelinks]{hyperref}
\usepackage{bookmark}

% ---- Enumerate ----
\usepackage{enumitem}

% ---- Diagramas ----
\usepackage{pgfplots}
\pgfplotsset{compat=1.18}
\usetikzlibrary{babel}  % babel español activa `>`; esto evita romper las flechas `->`
% courses/MetNum2023/Clases/assets/tikzstyles.tex
% ---------------------------------------------------------------------------
% Estilos compartidos para los diagramas del curso Métodos Numéricos
% (MetNum2023). Fuente única del "look" visual (colores, grosores, escalas)
% para que todas las figuras (matrices triangulares, curvas de convergencia
% / error, splines a trozos) sean consistentes entre clases.
%
% Uso: la clase debe cargar antes `tikz` y `pgfplots` (bloque §6.3 del
%      CLAUDE.md del curso), y luego
%      % courses/MetNum2023/Clases/assets/tikzstyles.tex
% ---------------------------------------------------------------------------
% Estilos compartidos para los diagramas del curso Métodos Numéricos
% (MetNum2023). Fuente única del "look" visual (colores, grosores, escalas)
% para que todas las figuras (matrices triangulares, curvas de convergencia
% / error, splines a trozos) sean consistentes entre clases.
%
% Uso: la clase debe cargar antes `tikz` y `pgfplots` (bloque §6.3 del
%      CLAUDE.md del curso), y luego
%      % courses/MetNum2023/Clases/assets/tikzstyles.tex
% ---------------------------------------------------------------------------
% Estilos compartidos para los diagramas del curso Métodos Numéricos
% (MetNum2023). Fuente única del "look" visual (colores, grosores, escalas)
% para que todas las figuras (matrices triangulares, curvas de convergencia
% / error, splines a trozos) sean consistentes entre clases.
%
% Uso: la clase debe cargar antes `tikz` y `pgfplots` (bloque §6.3 del
%      CLAUDE.md del curso), y luego
%      \input{../assets/tikzstyles.tex}
% (compilar desde el directorio de la clase para que la ruta relativa resuelva).
%
% Paleta propia de este curso (no se reusa la de Fisica3-2015/CDIV2017/
% ElecMag2024): mismos tonos base para alinear con keybox/notebox/warnbox,
% pero archivo independiente, sin dependencia cruzada entre cursos.
% ---------------------------------------------------------------------------

% ---- Paleta (alineada con las cajas del preámbulo: keybox/notebox/warnbox) ----
\definecolor{figblue}{HTML}{1F4E79}   % azul principal (L, curvas principales, keybox)
\definecolor{figred}{HTML}{B03A2E}    % rojo de énfasis (U, warnbox, error)
\definecolor{figamber}{HTML}{B8860B}  % ámbar (notebox, pivotes)
\definecolor{figgray}{HTML}{555555}   % auxiliares, ejes, ceros de matriz

% ---- TikZ: librerías y estilos reutilizables ----
% Nota: las puntas se declaran como `-{Latex[...]}` y no como `->`. La forma
% con llaves NO contiene un `>` literal, así que es inmune al `>` activo que
% introduce babel español (de todos modos se carga \usetikzlibrary{babel} en
% el bloque §6.3 para cubrir cualquier `->` suelto).
\usetikzlibrary{arrows.meta,calc,matrix,patterns}

\tikzset{
  figline/.style={figblue, line width=0.9pt},
  figaux/.style={figgray, dashed, line width=0.6pt},
  figaxis/.style={figgray, line width=0.7pt, -{Latex[length=2.2mm,width=1.5mm]}},
  % estructura de matrices (L triangular inferior, U triangular superior)
  matcelda/.style={minimum width=7mm, minimum height=7mm, anchor=center, font=\small},
  matnz/.style={matcelda, fill=figblue!12},
  matnzU/.style={matcelda, fill=figred!12},
  matcero/.style={matcelda, text=figgray!70},
  matpivote/.style={matcelda, fill=figamber!25, draw=figamber, line width=0.8pt},
  % nodos de interpolación / splines
  fignodo/.style={circle, inner sep=0pt, minimum size=4.5pt, draw=figblue, line width=0.9pt, fill=white},
  fignodoactivo/.style={fignodo, fill=figamber},
  % etiquetas
  figlbl/.style={font=\small},
  figlblaux/.style={font=\footnotesize, figgray},
}

% ---- pgfplots: estilos base ----
\pgfplotsset{
  % convergencia / error vs. n (grado, iteración) — suele ir en escala log
  errorfig/.style={
    width=10cm, height=6.5cm,
    xlabel style={font=\small},
    ylabel style={font=\small},
    tick label style={font=\footnotesize},
    every axis plot/.append style={line width=1.1pt, mark=*, mark size=1.5pt},
    grid=major, grid style={figgray!15},
    legend cell align=left,
    legend style={font=\footnotesize, draw=figgray!40, at={(0.5,-0.25)}, anchor=north},
    clip=false,
  },
  % interpolación / splines: un color por tramo, nodos marcados aparte
  interpfig/.style={
    width=10cm, height=6cm,
    axis lines=middle,
    xlabel style={font=\small, at={(axis description cs:1,0.5)}, anchor=west},
    ylabel style={font=\small, at={(axis description cs:0.5,1)}, anchor=south},
    tick label style={font=\footnotesize},
    every axis plot/.append style={line width=1.1pt},
    grid=none,
    legend cell align=left,
    legend style={font=\footnotesize, draw=figgray!40},
    clip=false,
  },
}

% (compilar desde el directorio de la clase para que la ruta relativa resuelva).
%
% Paleta propia de este curso (no se reusa la de Fisica3-2015/CDIV2017/
% ElecMag2024): mismos tonos base para alinear con keybox/notebox/warnbox,
% pero archivo independiente, sin dependencia cruzada entre cursos.
% ---------------------------------------------------------------------------

% ---- Paleta (alineada con las cajas del preámbulo: keybox/notebox/warnbox) ----
\definecolor{figblue}{HTML}{1F4E79}   % azul principal (L, curvas principales, keybox)
\definecolor{figred}{HTML}{B03A2E}    % rojo de énfasis (U, warnbox, error)
\definecolor{figamber}{HTML}{B8860B}  % ámbar (notebox, pivotes)
\definecolor{figgray}{HTML}{555555}   % auxiliares, ejes, ceros de matriz

% ---- TikZ: librerías y estilos reutilizables ----
% Nota: las puntas se declaran como `-{Latex[...]}` y no como `->`. La forma
% con llaves NO contiene un `>` literal, así que es inmune al `>` activo que
% introduce babel español (de todos modos se carga \usetikzlibrary{babel} en
% el bloque §6.3 para cubrir cualquier `->` suelto).
\usetikzlibrary{arrows.meta,calc,matrix,patterns}

\tikzset{
  figline/.style={figblue, line width=0.9pt},
  figaux/.style={figgray, dashed, line width=0.6pt},
  figaxis/.style={figgray, line width=0.7pt, -{Latex[length=2.2mm,width=1.5mm]}},
  % estructura de matrices (L triangular inferior, U triangular superior)
  matcelda/.style={minimum width=7mm, minimum height=7mm, anchor=center, font=\small},
  matnz/.style={matcelda, fill=figblue!12},
  matnzU/.style={matcelda, fill=figred!12},
  matcero/.style={matcelda, text=figgray!70},
  matpivote/.style={matcelda, fill=figamber!25, draw=figamber, line width=0.8pt},
  % nodos de interpolación / splines
  fignodo/.style={circle, inner sep=0pt, minimum size=4.5pt, draw=figblue, line width=0.9pt, fill=white},
  fignodoactivo/.style={fignodo, fill=figamber},
  % etiquetas
  figlbl/.style={font=\small},
  figlblaux/.style={font=\footnotesize, figgray},
}

% ---- pgfplots: estilos base ----
\pgfplotsset{
  % convergencia / error vs. n (grado, iteración) — suele ir en escala log
  errorfig/.style={
    width=10cm, height=6.5cm,
    xlabel style={font=\small},
    ylabel style={font=\small},
    tick label style={font=\footnotesize},
    every axis plot/.append style={line width=1.1pt, mark=*, mark size=1.5pt},
    grid=major, grid style={figgray!15},
    legend cell align=left,
    legend style={font=\footnotesize, draw=figgray!40, at={(0.5,-0.25)}, anchor=north},
    clip=false,
  },
  % interpolación / splines: un color por tramo, nodos marcados aparte
  interpfig/.style={
    width=10cm, height=6cm,
    axis lines=middle,
    xlabel style={font=\small, at={(axis description cs:1,0.5)}, anchor=west},
    ylabel style={font=\small, at={(axis description cs:0.5,1)}, anchor=south},
    tick label style={font=\footnotesize},
    every axis plot/.append style={line width=1.1pt},
    grid=none,
    legend cell align=left,
    legend style={font=\footnotesize, draw=figgray!40},
    clip=false,
  },
}

% (compilar desde el directorio de la clase para que la ruta relativa resuelva).
%
% Paleta propia de este curso (no se reusa la de Fisica3-2015/CDIV2017/
% ElecMag2024): mismos tonos base para alinear con keybox/notebox/warnbox,
% pero archivo independiente, sin dependencia cruzada entre cursos.
% ---------------------------------------------------------------------------

% ---- Paleta (alineada con las cajas del preámbulo: keybox/notebox/warnbox) ----
\definecolor{figblue}{HTML}{1F4E79}   % azul principal (L, curvas principales, keybox)
\definecolor{figred}{HTML}{B03A2E}    % rojo de énfasis (U, warnbox, error)
\definecolor{figamber}{HTML}{B8860B}  % ámbar (notebox, pivotes)
\definecolor{figgray}{HTML}{555555}   % auxiliares, ejes, ceros de matriz

% ---- TikZ: librerías y estilos reutilizables ----
% Nota: las puntas se declaran como `-{Latex[...]}` y no como `->`. La forma
% con llaves NO contiene un `>` literal, así que es inmune al `>` activo que
% introduce babel español (de todos modos se carga \usetikzlibrary{babel} en
% el bloque §6.3 para cubrir cualquier `->` suelto).
\usetikzlibrary{arrows.meta,calc,matrix,patterns}

\tikzset{
  figline/.style={figblue, line width=0.9pt},
  figaux/.style={figgray, dashed, line width=0.6pt},
  figaxis/.style={figgray, line width=0.7pt, -{Latex[length=2.2mm,width=1.5mm]}},
  % estructura de matrices (L triangular inferior, U triangular superior)
  matcelda/.style={minimum width=7mm, minimum height=7mm, anchor=center, font=\small},
  matnz/.style={matcelda, fill=figblue!12},
  matnzU/.style={matcelda, fill=figred!12},
  matcero/.style={matcelda, text=figgray!70},
  matpivote/.style={matcelda, fill=figamber!25, draw=figamber, line width=0.8pt},
  % nodos de interpolación / splines
  fignodo/.style={circle, inner sep=0pt, minimum size=4.5pt, draw=figblue, line width=0.9pt, fill=white},
  fignodoactivo/.style={fignodo, fill=figamber},
  % etiquetas
  figlbl/.style={font=\small},
  figlblaux/.style={font=\footnotesize, figgray},
}

% ---- pgfplots: estilos base ----
\pgfplotsset{
  % convergencia / error vs. n (grado, iteración) — suele ir en escala log
  errorfig/.style={
    width=10cm, height=6.5cm,
    xlabel style={font=\small},
    ylabel style={font=\small},
    tick label style={font=\footnotesize},
    every axis plot/.append style={line width=1.1pt, mark=*, mark size=1.5pt},
    grid=major, grid style={figgray!15},
    legend cell align=left,
    legend style={font=\footnotesize, draw=figgray!40, at={(0.5,-0.25)}, anchor=north},
    clip=false,
  },
  % interpolación / splines: un color por tramo, nodos marcados aparte
  interpfig/.style={
    width=10cm, height=6cm,
    axis lines=middle,
    xlabel style={font=\small, at={(axis description cs:1,0.5)}, anchor=west},
    ylabel style={font=\small, at={(axis description cs:0.5,1)}, anchor=south},
    tick label style={font=\footnotesize},
    every axis plot/.append style={line width=1.1pt},
    grid=none,
    legend cell align=left,
    legend style={font=\footnotesize, draw=figgray!40},
    clip=false,
  },
}


% ---- Miscellaneous ----
\usepackage{icomma}
\usepackage{siunitx}
\usepackage{textcomp}
\usepackage[bottom]{footmisc}

% ---- Title ----
\title{\textbf{Clase 12: Error de Interpolación} \\[0.3em]
        \large{Lagrange, Newton, diferencias divididas, y el teorema del error}}
\author{Transcripto y expandido --- Curso Métodos Numéricos (OpenFING)}
\date{}

\begin{document}

\maketitle
\thispagestyle{empty}
\newpage

\tableofcontents
\newpage

\section{La forma de Lagrange}

Segunda forma de escribir el polinomio interpolante, con una base
\textbf{adaptada a los nodos} $x_0,\dots,x_n$.

\textbf{Polinomio base de Lagrange} $L_n^k$: el único polinomio de grado
$\leq n$ tal que $L_n^k(x_j) = \delta_{jk}$ (delta de Kronecker). Fórmula
explícita:
$$L_n^k(x) = \prod_{\substack{i=0\\ i\neq k}}^{n} \frac{x - x_i}{x_k - x_i}$$

\begin{notebox}[Por qué el denominador nunca es cero]
Exactamente porque se exige $x_i$ distintos dos a dos. El numerador se
anula en todos los nodos salvo $x_k$; el denominador normaliza para que
valga $1$ en $x_k$.
\end{notebox}

\begin{keybox}[Polinomio interpolante en base de Lagrange]
$$p_n(x) = \sum_{k=0}^{n} y_k\, L_n^k(x)$$
\end{keybox}

\subsection{Ejemplo}

Datos: $(1/4,1),\ (1,-1),\ (5/2,3/2)$.
$$L_2^0(x) = \frac{(x-1)(x-5/2)}{(1/4-1)(1/4-5/2)}$$
(denominador $27/16$); análogamente $L_2^1,L_2^2$. El polinomio final es
el \textbf{mismo} que con Vandermonde, escrito en otra base.

\subsection{Luces y sombras de Lagrange}

\textbf{A favor}:
\begin{itemize}
  \item Ningún sistema que resolver $\to$ estable, sin problema de
  condicionamiento.
  \item Corregir un dato es fácil: cambia un número en la combinación
  lineal.
\end{itemize}

\textbf{En contra}:
\begin{itemize}
  \item Agregar un nodo obliga a recalcular \textbf{toda} la base (cada
  $L_n^k$ depende de todos los puntos).
  \item Evaluar en un punto que no es nodo es caro: hay que computar cada
  $L_n^k(x)$ por separado.
\end{itemize}

\begin{warnbox}
Esto encarece cualquier operación posterior (derivar, integrar, hallar
raíces), porque requieren evaluar el polinomio repetidamente.
\end{warnbox}

\section{La forma de Newton}

Recupera la estructura de la demostración por inducción (Clase 11):
$$p_n(x) = a_0 + a_1(x-x_0) + a_2(x-x_0)(x-x_1) + \dots + a_n\prod_{i=0}^{n-1}(x-x_i)$$

\subsection{Diferencias divididas: cómo calcular los coeficientes}

Notación $f[x_0,\dots,x_k]$, recursiva:
\begin{itemize}
  \item Orden 0: $f[x_i] = y_i$.
  \item Orden $k$: $f[x_i,\dots,x_{i+k}] = \dfrac{f[x_{i+1},\dots,x_{i+k}] - f[x_i,\dots,x_{i+k-1}]}{x_{i+k}-x_i}$.
\end{itemize}

Coeficientes de Newton: $a_k = f[x_0,\dots,x_k]$.

\textbf{Ejemplo} (mismos 3 puntos): $f[x_0]=-3/4$, $f[x_1]=-1$,
$f[x_2]=3/2$.
$$f[x_0,x_1] = \frac{-1-(-3/4)}{1-1/4} = -\frac13, \qquad f[x_1,x_2] = \frac{3/2-(-1)}{5/2-1} = \frac53$$
$$f[x_0,x_1,x_2] = \frac{5/3-(-1/3)}{5/2-1/4} = \frac{8}{9}$$
$$a_0=-\frac34,\quad a_1=-\frac13,\quad a_2=\frac89 \implies p_2(x) = -\frac34 -\frac13\Big(x-\frac14\Big) + \frac89\Big(x-\frac14\Big)(x-1)$$

\begin{notebox}[Estructura de árbol]
Se combinan pares consecutivos, luego esos resultados entre sí — una
``pirámide'' que sólo usa diferencias ya calculadas.
\end{notebox}

\subsection{Ventajas de Newton}

\begin{itemize}
  \item Ningún sistema que resolver $\to$ estable, igual que Lagrange.
  \item \textbf{Incremental}: agregar un punto sólo agrega un término
  $a_{n+1}\prod_{i=0}^{n}(x-x_i)$ — los coeficientes anteriores no
  cambian, porque $f[x_0,\dots,x_k]$ con $k<n+1$ no depende de $x_{n+1}$.
  Resuelve \textbf{ambos} problemas de Vandermonde a la vez.
  \item Evaluación eficiente vía el \textbf{algoritmo de Horner}
  (mencionado, no desarrollado en clase).
\end{itemize}

\section{Interpolando funciones desconocidas}

Nuevo problema: existe $f$ \textbf{desconocida}, sólo se tienen muestras
$(x_i,f(x_i))$. Pregunta central: ¿qué tan lejos está $p_n$ de $f$?

\subsection{La norma del supremo}

$$\|\varphi\|_{\infty,I} = \sup_{x\in I} |\varphi(x)|$$

(convergencia en esta norma = \textbf{convergencia uniforme}). Si $I$ es
compacto y $\varphi$ continua, por Weierstrass el supremo es un
\textbf{máximo}.

\section{El teorema de error de interpolación}

\begin{keybox}[Teorema]
Sea $f\in C^{n+1}([a,b])$, $x_0,\dots,x_n\in[a,b]$ distintos, $p_n$ el
polinomio interpolante. Para todo $x\in[a,b]$ existe $\gamma_x\in(a,b)$ tal
que
$$f(x) - p_n(x) = \frac{f^{(n+1)}(\gamma_x)}{(n+1)!}\,\omega_n(x), \qquad \omega_n(x) := \prod_{i=0}^{n}(x-x_i)$$
($\omega_n$ = \textbf{polinomio nodal}, grado $n+1$.)
\end{keybox}

\subsection{Demostración}

Fijado $x\neq x_i$ para todo $i$ (si $x=x_i$ ambos lados son $0$
trivialmente). Función auxiliar en la variable $t$:
$$g(t) = e_n(t) - \frac{e_n(x)}{\omega_n(x)}\,\omega_n(t), \qquad e_n:=f-p_n$$

\textbf{Raíces de $g$}: $g(x_i)=0$ para cada nodo (ambos términos se
anulan) y $g(x)=0$ (los $\omega_n$ se cancelan). $g$ tiene al menos $n+2$
raíces distintas en $[a,b]$.

\textbf{Rolle, repetido}: entre cada par de raíces consecutivas de $g$ hay
una raíz de $g'$ (Rolle): $g'$ tiene al menos $n+1$ raíces. Repitiendo:
$g''$ tiene al menos $n$, ..., $g^{(n+1)}$ tiene al menos $1$ raíz —
llamada $\gamma_x\in(a,b)$.

\begin{figure}[H]
\centering
\begin{tikzpicture}[scale=1.0]
  \node[figblue] at (4,1.3) {\footnotesize raíces de $g$};
  \draw[figaxis] (-0.3,0) -- (10.3,0);
  \foreach \x/\lbl in {0/x_0, 2/x_1, 4/x, 6/x_2, 8/x_3} {
    \filldraw[figblue] (\x,0) circle (2pt);
    \node[above, figblue] at (\x,0.2) {\footnotesize $\lbl$};
  }
  \foreach \x in {1,3,5,7} {
    \filldraw[figred] (\x,0) circle (1.8pt);
  }
  \node[below, figred] at (1,-0.3) {\footnotesize $z_0$};
  \node[below, figred] at (3,-0.3) {\footnotesize $z_1$};
  \node[below, figred] at (5,-0.3) {\footnotesize $z_2$};
  \node[below, figred] at (7,-0.3) {\footnotesize $z_3$};
  \node[figred] at (4,-1.1) {\footnotesize raíces de $g'$ (Rolle entre cada par consecutivo)};
\end{tikzpicture}
\caption{Con 4 nodos más $x$: $g$ tiene 5 raíces (azul); Rolle da 4 raíces
de $g'$ (rojo), una entre cada par consecutivo. Repitiendo Rolle sobre
$g',g'',\dots$ se llega a una única raíz de $g^{(n+1)}$: $\gamma_x$.}
\end{figure}

\textbf{Evaluar en $\gamma_x$}: $p_n$ tiene grado $\leq n$, así que
$e_n^{(n+1)} = f^{(n+1)}$. El término constante sobrevive la derivación.
$\omega_n$ tiene coeficiente principal $1$ y grado $n+1$, así que
$\omega_n^{(n+1)}(t)\equiv(n+1)!$. Por lo tanto:
$$0 = g^{(n+1)}(\gamma_x) = f^{(n+1)}(\gamma_x) - \frac{e_n(x)}{\omega_n(x)}\,(n+1)!$$
Despejando $e_n(x)$ se obtiene la fórmula del teorema. $\blacksquare$

\section{Corolario: cota práctica del error}

Tomando valor absoluto y acotando $|f^{(n+1)}(\gamma_x)|$ por el máximo de
$|f^{(n+1)}|$:
$$|f(x)-p_n(x)| \leq \frac{\|f^{(n+1)}\|_{\infty,[a,b]}}{(n+1)!}\,|\omega_n(x)|, \qquad \forall x\in[a,b]$$

Y tomando máximo en ambos lados:

\begin{keybox}
$$\|f-p_n\|_{\infty,[a,b]} \leq \frac{\|f^{(n+1)}\|_{\infty,[a,b]}}{(n+1)!}\,\|\omega_n\|_{\infty,[a,b]}$$
\end{keybox}

\begin{notebox}[Tres ingredientes en competencia]
Qué tan grande es $f^{(n+1)}$, el factorial $(n+1)!$ (siempre ayuda), y qué
tan grande es el polinomio nodal.
\end{notebox}

\section{¿Más puntos es mejor? Ejemplo con \texorpdfstring{$\sin x\cos x$}{sin(x)cos(x)}}

No necesariamente: el factorial compite contra el crecimiento de
$\|f^{(n+1)}\|_\infty$ con $n$.

\textbf{Ejemplo}: $f(x)=\sin(x)\cos(x)$ en $I=[0,1]$. Las derivadas
sucesivas combinan $\sin$ y $\cos$ (acotados por $1$), así que
$\|f^{(n+1)}\|_\infty \leq 2$ para todo $n$ (cota uniforme). Con nodos en
$[0,1]$, $|x-x_i|\leq 1$, así que $\|\omega_n\|_\infty \leq 1$
\textbf{sin importar cómo se elijan los nodos}. Entonces:
$$\|f-p_n\|_{\infty,[0,1]} \leq \frac{2^{n+1}}{(n+1)!} \xrightarrow[n\to\infty]{} 0$$

\begin{notebox}[Verificación numérica]
Con 11 puntos equiespaciados, el error ronda $10^{-11}$. La cota teórica da
$\sim 10^{-5}$ (grado 10) y $\sim 10^{-14}$ (grado 20) — agregar puntos
``hace pedazos'' el error en esta función. \emph{Contraste con la función
de Runge de la clase siguiente, donde el mismo razonamiento falla.}
\end{notebox}

\end{document}
